
\documentclass[class=llncs, crop=false]{standalone}
\usepackage{unicode-math}
\usepackage{newcomputermodern}
\usepackage{microtype}
\usepackage{newunicodechar}
\newunicodechar{⤳}{\ensuremath{⤳}}
\newunicodechar{↪}{\ensuremath{↪}}
\newunicodechar{≔}{\ensuremath{≔}}
\newunicodechar{⊢}{\ensuremath{⊢}}

% \usepackage{amsthm}
\usepackage{xcolor-material}
\usepackage{etoolbox}

%% ── Draft toggle ─────────────────────────────────────────────────
%% Comment out the next line for submission.
% \newcommand{\darkmode}{}

\ifdefined\darkmode
	\usepackage{showframe}
	\renewcommand\ShowFrameLinethickness{0.2pt}
	\renewcommand*\ShowFrameColor{\color{MaterialGrey600}}
	\usepackage{pagecolor}
	\pagecolor{MaterialGrey900}
	\color{MaterialGrey100}
\fi

\usepackage{tikz}
\usetikzlibrary{fit}
\usepackage{hyperref}
\ifdefined\darkmode
	\hypersetup{colorlinks=true,
		linkcolor=MaterialLightBlue300,
		citecolor=MaterialTeal300,
		urlcolor=MaterialLightBlue200}
\else
	\hypersetup{colorlinks=true,
		linkcolor=eoComment,
		citecolor=eoComment,
		urlcolor=eoComment}
\fi
\usepackage{ebproof}
\usepackage{array}

\usepackage{jbw-boxfigure}
\usepackage{jbw-fix-hyperref-autoref}

\usepackage{listings}
\usepackage[most]{tcolorbox}
\tcbuselibrary{listings}

%% ── Syntax colours ───────────────────────────────────────────────
\ifdefined\darkmode
	% Dark mode
	\definecolor{eoKeyword}{HTML}{f06292}%       commands   (pink)
	\definecolor{eoAttr}{HTML}{64b5f6}%          attributes (light blue)
	\definecolor{eoBuiltin}{HTML}{4dd0e1}%       builtins   (cyan)
	\definecolor{eoConst}{HTML}{b39ddb}%         types      (light purple)
	\definecolor{eoComment}{HTML}{9e9e9e}%       comments   (grey)
	\definecolor{eoString}{HTML}{ffb74d}%        strings    (light amber)
	\definecolor{eoPunct}{HTML}{bdbdbd}%         punctuation (light grey)
	\definecolor{eoFgMath}{HTML}{c8c0c8}%        foreground (softer)
\else
	% Light mode
	\definecolor{eoKeyword}{HTML}{e14775}%       commands   (pink)
	\definecolor{eoAttr}{HTML}{4a7fb5}%          attributes (muted blue)
	\definecolor{eoBuiltin}{HTML}{1c8ca8}%       builtins   (teal-cyan)
	\definecolor{eoConst}{HTML}{7058be}%         types      (purple)
	\definecolor{eoComment}{HTML}{3b5998}%       comments   (dark blue)
	\definecolor{eoString}{HTML}{cc7a0a}%        strings    (amber)
	\definecolor{eoPunct}{HTML}{918c8e}%         punctuation (grey)
	\definecolor{eoFgMath}{HTML}{3a3338}%        foreground (softer)
\fi

% Eunoia language definition for listings
\lstdefinelanguage{eunoia}{
	comment=[l]{;},
	string=[b]",
	morecomment=[s]{|}{|},% quoted symbols |...|
	alsoletter={-:@.},
	keywords={declare-const, declare-parameterized-const, declare-consts,
			declare-rule, declare-fun, declare-sort, declare-datatype,
			define, define-fun, define-const, define-sort,
			program, include, assume, assume-push, step, step-pop,
			let, match, as, par, forall, exists, lambda},
	keywords=[2]{:conclusion, :premises, :premise-list, :args,
			:requires, :assumption, :rule, :signature, :type,
			:right-assoc, :right-assoc-nil, :left-assoc, :left-assoc-nil,
			:chainable, :pairwise, :binder, :arg-list,
			:implicit, :opaque, :list, :sorry},
	% eo:: builtins handled via literate (see style below)
	keywords=[4]{Type, Proof},
	sensitive=true,
}

\lstdefinestyle{eo}{
language=eunoia,
basicstyle=\small\ttfamily\color{eoFgMath},
keywordstyle=\color{eoKeyword}\bfseries,
keywordstyle=[2]\color{eoAttr},
keywordstyle=[4]\color{eoConst},
commentstyle=\color{eoComment}\itshape,
stringstyle=\color{eoString},
showstringspaces=false,
columns=flexible,
keepspaces=true,
frame=none,
numbers=left,
numberstyle=\tiny\ttfamily\color{lineNoColor},
numbersep=6pt,
literate={(}{{\textcolor{eoPunct}{(}}}1
{)}{{\textcolor{eoPunct}{)}}}1
{[}{{\textcolor{eoPunct}{[}}}1
{]}{{\textcolor{eoPunct}{]}}}1
{eo::}{{\textcolor{eoBuiltin}{eo}\textcolor{eoPunct}{::}}}4
{$}{{\textcolor{eoConst}{\$}}}1,
	aboveskip=0.5\medskipamount,
	belowskip=0.5\medskipamount,
	}

	%% ── tcolorbox backgrounds ────────────────────────────────────────
	\ifdefined\darkmode
		\definecolor{eoBoxBg}{HTML}{2e3040}
		\definecolor{eoBoxFr}{HTML}{424454}
		\definecolor{lpBoxBg}{HTML}{2e3040}
		\definecolor{lpBoxFr}{HTML}{424454}
	\else
		\definecolor{eoBoxBg}{HTML}{fbfaf9}
		\definecolor{eoBoxFr}{HTML}{c5c1bf}
		\definecolor{lpBoxBg}{HTML}{fbfaf9}
		\definecolor{lpBoxFr}{HTML}{c5c1bf}
		\definecolor{lineNoColor}{HTML}{bfb9ba}
	\fi

	% tcolorbox style for Eunoia code blocks
	\newtcblisting{eolisting}{
		enhanced,
		listing only,
		listing options={style=eo, aboveskip=0pt, belowskip=0pt},
		colback=eoBoxBg,
		colframe=eoBoxFr,
		arc=0pt,
		boxrule=0.4pt,
		borderline west={1pt}{0pt}{eoBuiltin},
		left=14pt, right=2pt, top=2pt, bottom=2pt,
		before skip=\medskipamount,
		after skip=\medskipamount,
	}
	\newtcbinputlisting{\eofile}[1]{
		enhanced,
		listing only,
		listing file={#1},
		listing options={style=eo, aboveskip=0pt, belowskip=0pt},
		colback=eoBoxBg,
		colframe=eoBoxFr,
		arc=0pt,
		boxrule=0.4pt,
		borderline west={1pt}{0pt}{eoBuiltin},
		left=14pt, right=2pt, top=2pt, bottom=2pt,
		before skip=\medskipamount,
		after skip=\medskipamount,
		before upper=\listboxsetup,
	}

	% LambdaPi language definition and style
	\lstdefinelanguage{lambdapi}{
		comment=[l]{//},
		alsoletter={_},
		alsoother={@:$},
		keywords={symbol, rule, with, assert,
				require, open, type},
		keywords=[2]{constant, sequential, opaque, injective},
		keywords=[3]{TYPE},
		sensitive=true,
		mathescape=false,
		texcl=false,
	}

	\lstdefinestyle{lp}{
	language=lambdapi,
	basicstyle=\small\ttfamily\color{eoFgMath},
	keywordstyle=\color{eoKeyword}\bfseries,
	keywordstyle=[2]\color{eoAttr},
	keywordstyle=[3]\color{eoConst},
	commentstyle=\color{eoComment}\itshape,
	showstringspaces=false,
	columns=flexible,
	keepspaces=true,
	frame=none,
	numbers=left,
	numberstyle=\tiny\ttfamily\color{lineNoColor},
	numbersep=6pt,
	literate=
		{eo.Proof}{{\textcolor{eoBuiltin}{eo}\textcolor{eoPunct}{.}\textcolor{eoConst}{Proof}}}8
	{eo.app}{{\textcolor{eoBuiltin}{eo}\textcolor{eoPunct}{.}app}}6
	{eo.nil}{{\textcolor{eoBuiltin}{eo}\textcolor{eoPunct}{.}nil}}6
	{eo.}{{\textcolor{eoBuiltin}{eo}\textcolor{eoPunct}{.}}}3
	{eo::}{{\textcolor{eoBuiltin}{eo}\textcolor{eoPunct}{::}}}4
	{(}{{\textcolor{eoPunct}{(}}}1
	{)}{{\textcolor{eoPunct}{)}}}1
	{[}{{\textcolor{eoPunct}{[}}}1
	{]}{{\textcolor{eoPunct}{]}}}1
	{;}{{\textcolor{eoPunct}{;}}}1
	{:}{{\textcolor{eoPunct}{:}}}1
	{\{|}{{\textcolor{eoPunct}{\{|}}}2
	{|\}}{{\textcolor{eoPunct}{|\}}}}2
	{$}{{\textcolor{eoConst}{\$}}}1
{φ}{{$\varphi$}}1,
aboveskip=0.5\medskipamount,
belowskip=0.5\medskipamount,
}

% tcolorbox style for LambdaPi code blocks
\newtcblisting{lplisting}{
	enhanced,
	listing only,
	listing options={style=lp, aboveskip=0pt, belowskip=0pt},
	colback=lpBoxBg,
	colframe=lpBoxFr,
	arc=0pt,
	boxrule=0.4pt,
	borderline west={1pt}{0pt}{eoConst},
	left=14pt, right=2pt, top=2pt, bottom=2pt,
	before skip=\medskipamount,
	after skip=\medskipamount,
}
\newtcbinputlisting{\lpfile}[1]{
	enhanced,
	listing only,
	listing file={#1},
	listing options={style=lp, aboveskip=0pt, belowskip=0pt},
	colback=lpBoxBg,
	colframe=lpBoxFr,
	arc=0pt,
	boxrule=0.4pt,
	borderline west={1pt}{0pt}{eoConst},
	left=14pt, right=2pt, top=2pt, bottom=2pt,
	before skip=\medskipamount,
	after skip=\medskipamount,
	before upper=\listboxsetup,
}
% Thin separator line for use inside multi-part listing boxes
\newcommand{\listsep}{%
	\vspace{-\baselineskip}\vspace{7pt}%
	\noindent\kern-14pt\textcolor{eoBoxFr}{\rule{\dimexpr\linewidth+14pt+2pt}{0.3pt}}\par%
	\vspace{2pt}%
}
% Helper: zero out trivlist glue that listings inserts
\newcommand{\listboxsetup}{\topsep=0pt \partopsep=0pt \parsep=0pt \parskip=0pt}
% Multi-part listing boxes: use \listsep between \lstinputlisting calls
\newenvironment{eobox}{%
	\begin{tcolorbox}[
			enhanced, listing only,
			colback=eoBoxBg, colframe=eoBoxFr,
			arc=0pt, boxrule=0.4pt,
			borderline west={1pt}{0pt}{eoBuiltin},
			left=14pt, right=2pt, top=2pt, bottom=2pt,
			before skip=\medskipamount, after skip=\medskipamount,
		]\listboxsetup
		}{\end{tcolorbox}}
\newenvironment{lpbox}{%
	\begin{tcolorbox}[
			enhanced, listing only,
			colback=lpBoxBg, colframe=lpBoxFr,
			arc=0pt, boxrule=0.4pt,
			borderline west={1pt}{0pt}{eoConst},
			left=14pt, right=2pt, top=2pt, bottom=2pt,
			before skip=\medskipamount, after skip=\medskipamount,
		]\listboxsetup
		}{\end{tcolorbox}}
\lstset{defaultdialect=[]{eunoia}}

% Redefine example and definition to use italic headings
\let\example\relax
\let\endexample\relax
\spnewtheorem{example}{Example}{\itshape}{\rmfamily}
\def\exampleautorefname{example}
\let\definition\relax
\let\enddefinition\relax
\spnewtheorem{definition}{Definition}{\itshape}{\rmfamily}

\usepackage[backend=biber]{biblatex}
\addbibresource{inria.bib}

% small cases environment (brace + contents scaled to \small)
\newenvironment{scases}{%
	\left\{\,\small\begin{array}{@{}l@{\quad}l@{}}%
		}{%
	\end{array}\right.%
}

% font faces
\newcommand{\msf}{\mathsf}
\newcommand{\mbf}{\mathbf}
\newcommand{\mscr}{\mathscr}
\newcommand{\mtt}[1]{{\color{eoConst}{\texttt{#1}}}}% types/constants: Type, Bool, true, false, Proof
\newcommand{\kw}[1]{{\color{eoKeyword}{\texttt{\textbf{#1}}}}}% commands: declare-const, define, program, ...
\newcommand{\ev}[1]{{\color{eoFgMath}{\texttt{#1}}}}
\newcommand{\equo}{{\color{eoString}{\texttt{"}}}}% string quote
\newcommand{\estr}[1]{\equo{#1}\equo}% string literals

\newcommand{\cpc}{\msf{CPC}}
\newcommand{\tsm}[1]{{\small\text{#1}}}

\newcommand{\ptn}[1]{{\color{MaterialPurple500}{\texttt{#1}}}}
\newcommand{\mcomment}[1]{{\color{MaterialIndigo500}{\text{#1}}}}
\newcommand{\wrapp}[4]{{\color{#1}{#2}}{#4}{\color{#1}{#3}}}
\newcommand{\prn}[1]{\wrapp{MaterialGrey600}{\lparen}{\rparen}{#1}}

\newcommand{\cvc}[1]{\textsc{cvc{#1}}}
\newcommand{\verit}{\textsc{veriT}}
\newcommand{\eunoia}{\textsc{Eunoia}}

% ----- abstract syntax ----------
% plurality, binding
\newcommand{\plur}[2]{{#1}_{1}\ldots{#1}_{#2}}
\newcommand{\plurcomma}[2]{{#1}_{1}, \ldots, {#1}_{#2}}
\newcommand{\maybe}[1]{\wrapp{MaterialGrey600}{\langle}{\rangle_?}{#1}}
\newcommand{\many}[1]{\wrapp{MaterialGrey600}{\langle}{\rangle_+}{#1}}
\newcommand{\any}[1]{\wrapp{MaterialGrey600}{\langle}{\rangle_∗}{#1}}
\newcommand{\some}[1]{\wrapp{MaterialGrey600}{\langle}{\rangle}{#1}}
\newcommand{\binddot}[3]{{#1\,#2.\,#3}}

% binders
\newcommand{\lam}{\binddot{λ}}
\newcommand{\pii}{\binddot{Π}}
\newcommand{\alam}[1]{\binddot{\msf{λ}^{#1}}}
\newcommand{\apii}[1]{\binddot{\msf{Π}^{#1}}}
\newcommand{\lett}[2]{\msf{let}\ \prn{#1}\,\msf{in}\ #2}

% applications
\newcommand{\appSym}{\cdot}
\newcommand{\app}[2]{{#1}\appSym{#2}}
\newcommand{\apptwo}[3]{#1\,#2\,#3}
\newcommand{\appthree}[4]{#1\,#2\,#3\,#4}
\newcommand{\appfour}[5]{#1\,#2\,#3\,#4\,#5}
\newcommand{\appldots}[3]{\prn{\prn{\app{#1}{#2}}\,\ldots\,⋅\,{#3}}}

% universes
\newcommand{\PROP}{\mtt{PROP}}
\newcommand{\TYPE}{\mtt{TYPE}}
\newcommand{\KIND}{\mtt{KIND}}



% -------------- Eunoia ------------------------------------
% eunoia symbols
\newcommand{\epar}[1]{\wrapp{eoPunct}{\texttt{(}}{\texttt{)}}{#1}}
\newcommand{\eapp}[2]{\epar{#1\ #2}}
\newcommand{\earr}[1]{\epar{\ev{->}\ #1}}
\newcommand{\eprf}[1]{\epar{\mtt{Proof}\ #1}}
\newcommand{\ecase}[2]{\epar{#1\ #2}}
\newcommand{\eo}[1]{{\color{eoBuiltin}{\texttt{eo}}}{\color{eoPunct}{∷}}{\ev{#1}}}% builtins
\newcommand{\eotype}{\mtt{Type}}
\newcommand{\eoreqs}[3]{\eapp{\eo{requires}}{#1\ #2\ #3}}
\newcommand{\eoquote}[1]{\eapp{\eo{quote}}{#1}}

\newcommand{\attr}[1]{{\color{eoAttr}{\texttt{:#1}}}}% attributes
\newcommand{\premises}[1]{\attr{premises}\,\epar{#1}}
\newcommand{\premiselist}[2]{\attr{premise-list}\ #1\ #2}
\newcommand{\requires}[1]{\attr{requires}\ \epar{#1}}
\newcommand{\conclusion}[1]{
	\attr{conclusion}\ #1
}


\newcommand{\sym}[1]{\msf{s}_{\msf{#1}}}
% \newcommand{\eoapp}[2]{\prn{#1\ #2}}
% \newcommand{\earr}[1]{\mtt{#1}}
\newcommand{\ctyp}{\mbf{ctyp}}
\newcommand{\att}{\mbf{att}}

\newcommand{\prf}{\mbf{prf}}
\newcommand{\step}{\mbf{step}}

\newcommand{\eor}{\texttt{or}}
\newcommand{\econcat}[3]{\eapp{\eo{list\_concat}}{#1\ #2\ #3}}
\newcommand{\econs}[3]{\eapp{#1}{#2\ #3}}
\newcommand{\false}{\texttt{false}}

\newcommand{\rcn}[1]{\attr{right-assoc-nil}\ {#1}}
\newcommand{\rc}{\attr{right-assoc}}

\newcommand{\lcn}[1]{\attr{left-assoc-nil}\ {#1}}
\newcommand{\lc}{\attr{left-assoc}}

\newcommand{\ch}[1]{\attr{chainable}\ {#1}}
\newcommand{\pw}[1]{\attr{pairwise}\ {#1}}
\newcommand{\bd}[1]{\attr{binder}\ {#1}}
\newcommand{\al}[1]{\attr{arg-list}\ {#1}}
\newcommand{\eovar}[2]{\eapp{\eo{var}}{#1\ #2}}

% --- builtin operators (sec:eo-builtin) ---
\newcommand{\eoBuiltinFig}{
	\begin{array}[t]{r@{\ }l@{\ }l@{\quad}l}
		β & {⦂}
		  & \eo{ite}\ ∣\ \eo{eq}\ ∣\ \eo{requires}\ ∣\ \eo{typeof}\ ∣\ \eo{var}\ ∣\ \ldots
		  & \mcomment{(core)}
		\\
		  & {∣}
		  & \eo{and}\ ∣\ \eo{or}\ ∣\ \eo{not}\ ∣\ \eo{xor}
		  & \mcomment{(boolean)}
		\\
		  & {∣}
		  & \eo{add}\ ∣\ \eo{mul}\ ∣\ \eo{neg}\ ∣\ \eo{zdiv}\ ∣\ \ldots
		  & \mcomment{(arithmetic)}
		\\
		  & {∣}
		  & \eo{nil}\ ∣\ \eo{cons}\ ∣\ \eo{list\_concat}\ ∣\ \eo{list\_len}\ ∣\ \ldots
		  & \mcomment{(list)}
	\end{array}
}


\newcommand{\cons}{{\text{cons}}}
\newcommand{\nil}{{\text{nil}}}
\newcommand{\foldl}{\mbf{foldl}}
\newcommand{\foldr}{\mbf{foldr}}
\newcommand{\elab}[1]{\mbf{elab}_{\prn{#1}}}
\newcommand{\glue}[1]{\mbf{glue}_{\prn{#1}}}
\newcommand{\bool}{\texttt{Bool}}
% ---- eo commands ------
\newcommand{\dc}[2]{\epar{\kw{declare-const}\ {#1}\ {#2}}}
\newcommand{\dca}[3]{\epar{\kw{declare-const}\ {#1}\ {#2}\ {#3}}}

% --- figures for `eoas.tex` ------------------------------
% s & {⦂}\ \sym 0 ∣ \sym 1 ∣ \ldots
% 		& \mcomment{(symbols)}         &
%
\newcommand{\eoTermSyntax}{
	\begin{array}[t]{r@{\ }l@{\quad}l@{\qquad}r@{\ }l@{\quad}l}
		%!----------------------------------------!%
		t & {⦂}\ s \mid ℓ \mid \eapp{s}{\vec t} \mid \eapp{s}{\epar{\vec{ρ}}\ t} & \mcomment{(terms)}      &
		ρ & {⦂}\ \epar{s\ t\ \maybe{ν}}                                          & \mcomment{(parameters)}   \\[1mm]
		ℓ & {⦂}\ \msf{int}(n) ∣ \msf{str}(s)                                     & \mcomment{(literals)}   &
		c & {⦂}\ \epar{t\ t'}                                                    & \mcomment{(cases)}        \\[1mm]
	\end{array}
}

\newcommand{\eoAttrSyntax}{
	\begin{array}[t]{r@{\ }c@{\ }l@{\quad}l}
		ν & {⦂} & \attr{implicit} \mid \attr{list}               & \mcomment{(prm. attributes)}   \\
		α & {⦂} & \attr{right-assoc} ∣ \attr{right-assoc-nil}\ t & \mcomment{(const. attributes)} \\
		  & {∣} & \attr{left-assoc} ∣ \attr{left-assoc-nil}\ t                                    \\
		  & {∣} & \attr{chainable}\ s ∣ \attr{pairwise}\ s                                        \\
		  & {∣} & \attr{binder}\ s ∣ \attr{arg-list}\ s
	\end{array}
}

\newcommand{\eoCommandSyntax}{
	\begin{array}[t]{r@{\ }l@{\quad}l}
		δ & \begin{array}[t]{c@{\ }l}
			    {⦂} & \epar{\kw{declare-const}\ s\ t\ \maybe{α}}                               \\
			    {∣} & \epar{\kw{declare-parameterized-const}\ s\ \epar{\vec{ρ}}\ t\ \maybe{α}} \\
			    {∣} & \epar{\kw{declare-consts}\ ℓ\ t}                                         \\
			    {∣} & \epar{\kw{define}\ {s} \ \epar{\vec ρ} \ t\ \ \maybe{\attr{type}\,t'}}   \\
			    {∣} & \epar{\kw{program}\ {s} \ \epar{\vec ρ}\
			    \attr{signature}\,\epar{\vec{t}}\,t'\
			    \epar{\vec{c}}
			    }                                                                              \\
			    {∣} & {\color{eoPunct}{\texttt{(}}}
			    \kw{declare-rule}\ s\ \epar{\vec ρ}\
			    \\ & \quad \maybe{\attr{assumption}\ t_{\text{assm}}}\
			    \\ & \quad \maybe{
			    \attr{premises}\ \epar{\vec{t}_{\text{prem}}}
			    \mid
			    \attr{premise-list}\ t_{\text{prem}}\ t_{\text{op}}
			    }\
			    \\ & \quad \maybe{\attr{args}\ \epar{\vec{t}_{\text{args}}}}\
			    \\ & \quad \maybe{\attr{requires}\ \epar{\vec{c}}}\
			    \\ & \quad \attr{conclusion}\ {t_{\text{conc}}}
			    {\color{eoPunct}{\texttt{)}}}                                                  \\
			    {∣} & \epar{\kw{include}\ μ}
		    \end{array}
		  & \mcomment{(std. commands)}
	\end{array}
}

\newcommand{\eoPrfCommandSyntax}{
	\begin{array}[t]{r@{\ }l@{\quad}l}
		π & \begin{array}[t]{c@{\ }l}
			    {⦂} & \epar{\kw{assume}\ s\ φ}      \\
			    {∣} & \epar{\kw{assume-push}\ s\ φ} \\
			    {∣} & \epar{\kw{step}\ s\ φ\
			    \attr{rule}\ s'\
			    \maybe{\attr{premises}\ {\vec ψ}}\
			    \maybe{\attr{args}\ {\vec t}}}      \\
			    {∣} & \epar{\kw{step-pop}\ s\ φ\
			    \attr{rule}\ s'\
			    \maybe{\attr{premises}\ {\vec ψ}}\
			    \maybe{\attr{args}\ {\vec t}}}      \\
		    \end{array}
		  & \mcomment{(prf. commands)}
	\end{array}
}

%
%
% --- macros for `lp.tex` ------------------------------
\newcommand{\esc}[1]{\wrapp{MaterialGrey600}{⦃}{⦄}{#1}}
\newcommand{\var}[1]{\msf{v}_{#1}}
\newcommand{\lbl}[2]{\rname{#1}\quad{#2}}
\newcommand{\lpTermSyntax}{
	\begin{array}[t]{r@{\ }l@{\quad}r}
		% ------------------------------------- %
		μ & {⦂}\ \TYPE ∣ \KIND
		  & \mcomment{(universes)}                                                             \\[1mm]
		% ------------------------------------- %
		t & {⦂}\ x ∣ κ ∣ μ ∣ \prn{\app{t}{t'}} ∣ \prn{\lam{x : t}{t'}} ∣ \prn{\pii{x : t}{t'}}
		  & \mcomment{(terms)}
		% ------------------------------------- %
	\end{array}
}

\newcommand{\lpCtxSyntax}{
	\begin{array}[t]{r@{\ }l@{\quad}r}
		% ------------------------------------- %
		γ & {⦂}\ \prn{x : t} \mid \prn{t ↪ t'} \
		  & \mcomment{(context elements)}        \\
		% ------------------------------------- %
		Γ & {⦂}\ \any{γ}
		  & \mcomment{(contexts)}
		% ------------------------------------- %
	\end{array}
}

\newcommand{\judge}[3]{{#1} ⊢_{Σ} {#2} : {#3}}
\newcommand{\wf}[1]{\msf{wf}\,{#1}}
% typing rules
\newcommand{\typeRule}{
	\begin{prooftree}
		\hypo{\wf Γ}
		\infer1{ \judge Γ \TYPE \KIND }
	\end{prooftree}
}

\newcommand{\varRule}{
	\begin{prooftree}
		\hypo{ \wf Γ}
		\infer1[$\prn{x : t}∈ Γ$]{ \judge Γ x t }
	\end{prooftree}
}

\newcommand{\constRule}{
	\begin{prooftree}
		\hypo{\wf Γ}
		\hypo{\judge {} t μ}
		\infer2[$\prn{κ : t}∈ Σ$]{ \judge Γ κ t }
	\end{prooftree}
}

\newcommand{\prodRule}{
	\begin{prooftree}
		\hypo{\judge Γ t {\TYPE}}
		\hypo{\judge{Γ, \prn{x : t}}{t'}{μ'}}
		\infer2{\judge Γ {\prn{\pii {x : t} {t'}}} {μ'} }
	\end{prooftree}
}

\newcommand{\absRule}{
	\begin{prooftree}
		\hypo{\judge{Γ, \prn{x : t}} e {t'}}
		\hypo{\judge Γ {\prn{\pii {x : t} {t'}}} μ}
		\infer2{\judge Γ {\prn{\lam {x : t} {e}}} {\prn{\pii {x : t} {t'}}} }
	\end{prooftree}
}

\newcommand{\appRule}{
	\begin{prooftree}
		\hypo{\judge Γ e {\prn{\pii{x : t}{t'}}}}
		\hypo{\judge Γ {e'} t}
		\infer2{\judge Γ {\prn{\app{e}{e'}}} {t'[x ↦ {e'}]}}
	\end{prooftree}
}

\newcommand{\convRule}{
	\begin{prooftree}
		\hypo{\judge Γ e t}
		\hypo{\judge Γ {t'} {μ}}
		\infer2[$\prn{t ≡_{βΣ} t'}$]{\judge Γ {e} {t'}}
	\end{prooftree}
}

\newcommand{\wfRuleEmp}{
	\begin{prooftree}
		\infer0{ \wf{\msf{∅}} }
	\end{prooftree}
}

\newcommand{\wfRuleVar}{
	\begin{prooftree}
		\hypo{ \judge Γ t μ }
		\infer1[$x ∉ \mbf{dom}(Γ)$]{ \wf{(Γ, \prn{x : t})}}
	\end{prooftree}
}

\newcommand{\rname}[1]{\mcomment{($\textsc{#1}$)}}

\newcommand{\lpTypingRules}{
	\begin{array}[t]{c}
		\lbl{wf0}{\wfRuleEmp} \qquad \lbl{wf+}{\wfRuleVar}
		\\[6mm]
		\lbl{var}{\varRule}    \qquad  \lbl{con}{\constRule}
		\\[6mm]
		\lbl{univ}{\typeRule}  \qquad \lbl{prod}{\prodRule}
		\\[6mm]
		\lbl{fun}{\absRule} \\[6mm]
		\lbl{app}{\appRule} \\[6mm]
		\lbl{conv}{\convRule}
	\end{array}
}



% --- concrete syntax for lambdapi
\newcommand{\lambdapi}{\textsc{LambdaPi}}
\newcommand{\llam}{\binddot{\msf{λ}}}
\newcommand{\ppii}{\binddot{\msf{Π}}}

% \newcommand{\symbol}{\mtt{symbol}}

\newcommand{\prm}[1]{\mbf{#1}}
\newcommand{\xsym}[1]{\mtt{@}{#1}}
\newcommand{\meta}[1]{\ptn{?{#1}}}

\newcommand{\ptrn}[1]{\ptn{\$#1}}
\newcommand{\pv}{\mbf{pvars}}

\newcommand{\imp}[1]{[#1]}
\newcommand{\xpl}[1]{\left(#1\right)}

\newcommand{\lpTermsConcrete}{
	\begin{array}[t]{r@{\ }l@{\quad}r}
		t & {⦂}\ s ∣ \imp{t} ∣ \prn{\app{t}{t'}}
		∣ \prn{\llam{ρ}{t}} ∣ \prn{\ppii{ρ}{t}}
		  & \mcomment{(terms)}
	\end{array}
}

\newcommand{\lpPatternSyntax}{
	\begin{array}[t]{r@{\ }c@{\ }l@{\quad}r@{\qquad}r@{\ }c@{\ }l@{\quad}r}
		m & {⦂}
		  & \mtt{constant} ∣ \mtt{sequential}
		  & \mcomment{(modifiers)}
		  &
		θ & {⦂}
		  & \ptrn{x} ∣ s\,\any{θ}
		  & \mcomment{(patterns)}
		\\[1mm]
		ρ & {⦂}
		  & \xpl{s : t} ∣ \imp{s : t}
		  & \mcomment{(parameters)}
		  &
		% ------------------------------------- %
		r & {⦂}
		  & \prn{{s\,\any{θ}} ↪ {θ'}}
		  & \mcomment{(rewrite rules)}
		% ------------------------------------- %
	\end{array}
}

\newcommand{\lpConcreteSyntax}{
	\begin{array}[t]{r@{\ }c@{\ }l@{\quad}r}
		c & {⦂}
		  & \maybe{m}\, \mtt{symbol}\ s \ \any{ρ} \,{: t} \ \maybe{≔ t'};
		  & \mcomment{(commands)}
		\\
		  & {∣}
		  & \mtt{rule}\ r\ \any{\mtt{with}\ r'};
		\\
		  & {∣}
		  & \mtt{require}\,\mtt{open}\,\many{μ};
	\end{array}}

% \newcommand{\symb}[1]{\mtt{symbol}\ {#1};}

% \newcommand{\symb}[1]{\mtt{rule}\ {#1};}

\newcommand{\Set}{\msf{Set}}
\newcommand{\El}[1]{\msf{τ}\,{#1}}

% ------- translation stuff -----------
\newcommand{\Term}{\msf{Term}}

%
\newcommand{\EO}{\msf{EO}}
\newcommand{\LP}{\msf{LP}}
\newcommand{\trans}[1]{\left⟦{#1}\right⟧}
\newcommand{\tm}[1]{\trans{#1}_\msf{tm}}
\newcommand{\ty}[1]{\trans{#1}_\msf{ty}}
\newcommand{\cmd}[1]{\trans{#1}_\msf{cmd}}
\newcommand{\ctx}[1]{\trans{#1}_\msf{ctx}}


\begin{document}
%
\newpage
\section{Encoding Eunoia in the λΠ-calculus}\label{sec:lambdapi}
\begin{boxfigure}[t!]{fig:lpTermSyntax}
	{Syntax and typing rules for the $λΠ$-calculus modulo rewriting.}
	\alttext{Top: term grammar for the lambda-Pi calculus modulo rewriting.
		Universes mu range over TYPE and KIND.
		Terms t range over variables, constants, universes, applications,
		lambda abstractions, and dependent products.
		Bottom: natural-deduction typing rules---empty context is well-formed
		(wf0); extending with (x:t) requires Gamma deriving t at some universe
		(wf+); variable rule reads (x:t) from the context (var);
		constant rule reads (kappa:t) from the signature with the constant's
		type sortable in Gamma (con); TYPE has type KIND (univ);
		dependent product, abstraction, application, and conversion modulo
		beta and rewriting (prod, fun, app, conv).}{%
		\vspace{4pt}\centerline{$\lpTermSyntax$}
		\boxfigurehline
		\centerline{$\lpTypingRules$}\vspace{4pt}}
\end{boxfigure}
%
\Autoref{fig:lpTermSyntax} provides syntax and typing rules
for the \emph{$λΠ$-calculus modulo rewriting}.
%
Each term is either a \emph{variable} $x$,
a \emph{constant} $κ$,
a \emph{universe} from $\{ \TYPE, \KIND \}$,
an \emph{application} of two terms $\prn{\app{t}{t'}}$,
or an \emph{abstraction}
$\prn{\binddot{\mscr{B}}{x : t}{t'}}$
where $\mscr{B}$ is a \emph{binder} from $\{ λ, Π \}$.
%
Terms are identified up to $α$-conversion
(i.e., renaming of bound variables), and we assume
appropriate definitions for \emph{substitution} $\prn{t[x ↦ t']}$
and \emph{$β$-reduction} $\prn{t ⇝_{β} t'}$.
%
Hereinafter, we may write $\prn{t_1 → t_2}$ for the term
$\prn{\ppii{(x : t_1)}{t_2}}$, where $x$ does not occur
free in $t_2$.

% --------------------------------------------------------%
A \emph{typing} is an expression of the form $\prn{t : t'}$,
and a \emph{context} is a list of typings each of the
form $\prn{x : t}$.
%
% Hereinafter, given a context $Γ$ and a typing $\prn{x : t}$,
% let $\prn{Γ, \prn{x : t}}$ denote the extension $Γ ∪ \{x : t\}$.
%
A \emph{rewrite rule} is an expression of the form
$\prn{t ↪ t'}$, where $t$ has the form
$\prn{\app{\prn{\app{κ}{t_1}}\ \ldots}{t_n}}$.
%
A \emph{signature} is a list of typings
$\prn{κ : t}$ (with $t$ closed) and rewrite rules.
%
For any signature $Σ$, let $R_{Σ}$ be the smallest
binary relation such that:
%
\begin{enumerate}
	\item for any rule $\prn{ℓ ↪ r} ∈ Σ$ and substitution $σ$,
	      $(ℓσ, rσ) ∈ R_{Σ}$, and
	\item $R_{Σ}$ is \emph{congruent} under application
	      and abstraction.
\end{enumerate}
%
%, abstraction and substitution.  $R$
Then, \emph{equality modulo rewriting} $\prn{{≡_{βΣ}}}$ is
defined as the least equivalence relation containing
$R_{Σ}$ and the $β$-reduction relation.
%
The rules in \autoref{fig:lpTermSyntax}
provide a (mutually inductive) definition of a
\emph{well-formedness} relation $\prn{\msf{wf}_Σ}$
on contexts and a \emph{typing relation}
$\prn{⊢_{Σ}}$ between contexts and typings,
where $\prn{Γ ⊢_{Σ} e : t}$ may be read as
``$e$ has type $t$ with respect to signature $Σ$ and context $Γ$''.

\begin{remark}
	The typing rules of the $λΠ$-calculus disallow
	abstractions on types (i.e., there are no
	expressions of type $\TYPE → t$ in any context).
	Eunoia's types are therefore encoded as $λΠ$-terms
	of type $\Set$, using the Tarski-style universe of
	LambdaPi's standard library: an injective symbol
	$\msf{τ} : \Set → \TYPE$ lifts a $\Set$-term to a
	$λΠ$-type, and the (infix) function-space symbol
	$\prn{⤳} : \Set → \Set → \Set$ satisfies
	$\msf{τ}\,\prn{A ⤳ B} ↪ \msf{τ}\,A → \msf{τ}\,B$.

	The dependent types arising from $\eo{quote}$
	arguments are not expressible with $\prn{⤳}$
	alone, so we additionally introduce a dependent
	function space $\prn{⤳^{\msf d}} :
		\pii{T : \Set}{\prn{{\msf{τ}\,T → \Set}} → \Set}$
	with the rule
	$\msf{τ}\,\prn{T ⤳^{\msf d} F} ↪
		\pii{x : \msf{τ}\,T}{\msf{τ}\,\prn{\app F x}}$;
	the non-dependent case is recovered by taking
	$F ≔ \lam{\_}{B}$.
\end{remark}

% -------------------------------------------------------- %
\subsection{Translating Eunoia Terms}
%
We begin by defining the translation $\tm{⋅}$ on
Eunoia terms; it produces $λΠ$-terms in the
abstract calculus of \autoref{fig:lpTermSyntax} and
assumes all terms have already been processed by the
$\mbf{elab}$ operator of \autoref{sec:eo-elab}.
%
Because some Eunoia symbols are not valid LambdaPi
identifiers, the translation on symbols uses LambdaPi
syntax for escaping: $\bar s ≔ \esc s$ iff $s$
contains forbidden characters, and $\bar s ≔ s$
otherwise.
%
\begin{definition}\label{def:tm-ty}
	%
	Let $\tm{⋅}$ be the least operator such that:

	\begin{enumerate}
		\item For any symbol $s$ or literal $ℓ$, the following hold:
		      $$\tm{s} =
			      \begin{scases}
				      \lpo{Type}   & \text{if $s = \eotype$,} \\
				      \bar s & \text{otherwise.}
			      \end{scases}
			      %
			      \quad
			      %
			      \tm{ℓ} =
			      \begin{scases}
				      \ev{of\_Z}\ \tm{n}
				      &
				      \text{if $ℓ = \msf{int}(n)$,}
				      \\
				      \ev{of\_Q}\ \tm{n}\ \tm{m}
				      &
				      \text{if $ℓ = \msf{rat}(n,m)$,}
				      \\
				      \bar s & \text{if $ℓ = \msf{str}(s)$.}
			      \end{scases}
		      $$
		      The symbols $\ev{of\_Z}$ and $\ev{of\_Q}$ are
		      defined in our prelude (\autoref{sec:results}) and embed
		      the integer and rational literals of LambdaPi's
		      standard library into the carrier type of Eunoia's
		      arithmetic sorts.
		\item For arrow types, the following hold,
		      where $\epar{x\ T} ∈ \vec ρ$:
		      $$
			      \begin{array}[t]{r@{\ {=}\ }l}
				      \tm{\earr{\eoquote{x}\ {\vec t}}}
				                    &
				      \tm{T} ⤳^{\msf d} \lam{x}{\tm{\earr{\vec t}}}
				      \\[1mm]
				      \tm{\earr{t'\ {\vec t}}}
				                    &
				      \tm{t'} ⤳ \tm{\earr{\vec t}}
				      \\[1mm]
				      \tm{\earr{t}} & \tm t
			      \end{array}
		      $$
		\item For binary and $n$-ary applications,
		      the following hold:
		      $$
			      \begin{array}[t]{l@{\ {=}\ }l}
				      \tm{\eapp{\mtt{\_}}{t_1\ t_2}}
				       & \lpo{app}\ \tm{t_1}\ \tm{t_2}
				      \\[2mm]
				      \tm{\eprf{t}}
				       & \lpo{Proof}\ {\tm {t}}
				      \\[2mm]
				      \tm{\eapp{s}{\vec t}}
				       &
				      \bar s\ \tm{t_1}\ \ldots\ \tm{t_n}
			      \end{array}
		      $$
		\item The only normal-form binder application is $\eo{define}$,
		      which is translated using LambdaPi's standard
		      $\msf{let}$ binder:
		      $$\tm{\eapp{\eo{define}}{
					      \epar{\vec v}
					      \ t'
				      }
			      }
			      = \prn{
				      \msf{let}\ \vec {w}
				      \ \msf{in}\ \tm{t'}
			      }$$
		      where $w_i = \prn{{\bar x}_i ≔ \tm{t_i}}$
		      for each $v_i = \epar{x_i\ t_i}$.
	\end{enumerate}
	%
\end{definition}

% -------------------------------------------------------- %
\subsection{The LambdaPi Proof Assistant}
%
\begin{boxfigure}[t!]{fig:lpSyntax}
	{Syntax for a fragment of the LambdaPi proof assistant.}
	\alttext{Two grammar tables for LambdaPi concrete syntax.
		Top: modifiers (constant, sequential); parameters (explicit (s:t) and
		implicit [s:t]); patterns theta (pattern variables \$x or pattern
		applications); rewrite rules (s applied to patterns rewriting to theta).
		Bottom: commands---symbol declaration with optional modifier, parameters,
		type, and optional definition; rule declaration listing rewrite rules with
		'with' separators; assertion 'assert |- t : t';
		and 'require open' for module imports.}{%
		\vspace{4pt}\centerline{$\lpPatternSyntax$}
		\boxfigurehline
		\centerline{$\lpConcreteSyntax$}\vspace{4pt}}
\end{boxfigure}
%
\Autoref{fig:lpSyntax} defines the fragment of LambdaPi's
concrete syntax that the translation produces. It is
sugar over the calculus of \autoref{fig:lpTermSyntax}:
implicit binders $\imp{x : t}$ and placeholders $\_$
are elaborated to $Π$-binders and metavariables during
type inference, and each declaration extends~$Σ$.

A \emph{symbol declaration}
$\maybe{m}\,\kw{symbol}\ s\ \vec ρ : t\,\maybe{≔ t'}$
introduces $s$ with the typing
$\prn{\xsym s : \pii{\vec ρ}{t}}$, where parameters
$\vec ρ$ are either \emph{explicit} $\xpl{x : t}$ or
\emph{implicit} $\imp{x : t}$.
We write $\xsym s$ (with an @) for the
\emph{fully-explicit} form; an unprefixed~$s$ inserts a
placeholder for each leading implicit, to be filled by
higher-order unification.
The optional $\md{constant}$ modifier forbids rewrite
rules with $s$ at the head; $\md{sequential}$ matches
rewrite rules with head $s$ in their textual order; and
a definition $\prn{≔ t'}$ adds the rewrite rule
$\prn{s ↪ \lam{\vec ρ}{t'}}$.
%
A \emph{rewrite rule declaration} $\kw{rule}\ θ ↪ θ'$
adds the rule $\prn{θ ↪ θ'}$ to the signature, where
$θ$ and $θ'$ are \emph{patterns}---either a
\emph{pattern variable} $\ptrn{x}$ (written \$x) or a
\emph{pattern application} $\prn{s\,\any{θ}}$.
The rule is valid when $θ$ is a pattern application and
every pattern variable in $θ'$ also occurs in $θ$.
LambdaPi additionally verifies \emph{subject reduction}
at declaration time: for every typed substitution of
the pattern variables, both sides must have the same
type. An \emph{assertion} $\kw{assert}\ ⊢\ t : t'$
checks $t : t'$ without extending the signature.

% ----------------------------------------------------------
\subsection{Translating Eunoia Commands and Proofs}
% ----------------------------------------------------------
We extend the term translation $\tm{⋅}$
(\autoref{def:tm-ty}) to parameter contexts and to
both standard and proof commands.
%
\Autoref{tab:encoding} summarises the full
Eunoia-to-LambdaPi encoding; the remainder of this
section formalises the operators on contexts
(\autoref{def:ctx}), standard commands
(\autoref{def:cmd}), and proof commands
(\autoref{def:cmd-prf}).

\begin{table}[t!]
	\centering
	\caption{Eunoia-to-LambdaPi encoding, grouped by the kind of design decision involved.
		Symbols prefixed with $\msf{eo.}$ live in our hand-written
		prelude (\autoref{sec:results}); $⤳$ and $⤳^{\msf d}$
		are the function-space symbols of \autoref{fig:lpTermSyntax}.}
	\label{tab:encoding}
	\small
	\begin{tabular}{@{}p{0.36\linewidth}@{\ }p{0.56\linewidth}@{}}
		\hline
		\textbf{Eunoia feature} & \textbf{LambdaPi encoding} \\
		\hline
		\multicolumn{2}{@{}l@{}}{\textit{Types via a Tarski universe (LambdaPi disallows abstractions over types)}}
		\\
		Type universe $\eotype$
		& sort $T$ becomes a term $T : \Set$ with carrier $\msf{τ}\,T$
		\\
		Function type $\earr{A\ B}$
		& $A ⤳ B$, with rewrite
		$\msf{τ}\,(A ⤳ B) ↪ \msf{τ}\,A → \msf{τ}\,B$
		\\
		Object-level dependent type from $\eoquote{x}$
		& $A ⤳^{\msf d} \lam{x}{B}$ (Remark~1; matchable head, unlike~$Π$)
		\\
		\hline
		\multicolumn{2}{@{}l@{}}{\textit{$n$-ary syntax compiled to binary applications via $\mbf{elab}$ (\autoref{def:elab})}}
		\\
		Higher-order application $\eho{t_1}{t_2}$
		& $\lpo{app}\ t_1\ t_2$
		\\
		$n$-ary $\eapp{f}{\vec t}$ with $\rcn{t_\nil}$, $\ch{}$ etc.\
		& nested binary application; nil terminator handled by a rewrite rule
		\\
		Variables in binder applications ($\eo{var}$)
		& reified into a heterogeneous-list type whose entries carry their own sort
		\\
		\hline
		\multicolumn{2}{@{}l@{}}{\textit{Programs and rules as computational/axiomatic content}}
		\\
		Program (\kw{program})
		& \md{sequential}\ \kw{symbol}\ + one rewrite rule per case (SR-checked)
		\\
		Inference rule (\kw{declare-rule})
		& \md{constant}\ \kw{symbol}\ + auxiliary symbol with quote-pattern rewrite rule when $\attr{args}$ is present
		\\
		Built-in arithmetic ($\eo{add}$, $\eo{mul}$, $\ldots$)
		& prelude rewrite rules delegating to LambdaPi's $\mathbb Z, \mathbb Q$
		\\
		Conclusion guard $\eoreqs{l}{r}{ψ}$
		& $ψ$ applies only when $l$ and $r$ are syntactically identical
		\\
		\hline
		\multicolumn{2}{@{}l@{}}{\textit{Straightforward}}
		\\
		Parameter attribute $\attr{implicit}$
		& implicit LambdaPi binder $\imp{x : \msf{τ}\,T}$
		\\
		Constant decl / macro \kw{define}
		& \kw{symbol}\ $s\,\vec ρ : \msf{τ}\,T$ resp.\ $\kw{symbol}\ s\,\vec ρ ≔ t$
		\\
		Proof commands \kw{assume} / \kw{step}
		& \md{constant}\ \kw{symbol} of type $\lpo{Proof}\,φ$, resp.\ a defined symbol followed by $\kw{assert}$
		\\
		\hline
	\end{tabular}
\end{table}


% ----------------------------------------------------------
\subsubsection{Translating Commands.}
% ----------------------------------------------------------
Parameter contexts and two auxiliary notations
support the command translation that follows.

\begin{definition}\label{def:ctx}
	Given $\vec ρ = \epar{x_1\ t_1\ \maybe{ν_1}}
		\ldots \epar{x_n\ t_n\ \maybe{ν_n}}$, define:
	$$
		\ctx{\vec ρ}\ ≔\
		ρ_1^★\, \ldots\, ρ_n^★
		\quad\text{where}\quad
		ρ_i^★ =
		\begin{scases}
			\imp{\bar x_i : \El{\tm{t_i}}}
			& \text{if $ν_i = \attr{implicit}$,} \\
			\xpl{\bar x_i : \El{\tm{t_i}}}
			& \text{otherwise.}
		\end{scases}
	$$
\end{definition}
%
We further write $\tm{t}^{\ptrn{}}$ for the result of
applying $\tm{⋅}$ to $t$ and replacing each free
variable $x$ with the pattern variable $\ptrn{x}$;
and $\mbf{impl}(\vec ρ, t)$ for the sublist of
$\ctx{\vec ρ}$ containing each $x_i ∈ \vec ρ$ that
occurs free in $t$, marked as implicit.
%
\begin{definition}\label{def:cmd}
	%
	Let $δ$ be a standard Eunoia command. Then $\cmd{δ}$
	is the list of LambdaPi commands described as follows:
	%
	\begin{enumerate}\itemsep3pt
		\item If $δ$ is a constant declaration,
		      $δ ⊢ s(\vec ρ) : t$, the head of
		      $\cmd δ$ is
		      $\kw{symbol}\ \bar s\ \ctx{\vec ρ}
			      : \El{\tm t}\mtt{;}$.
		      If $δ ⊢ s ∷ α$ and $α$ is
		      $\rcn{t_\nil}$ or $\lcn{t_\nil}$, the
		      following rewrite rule is also produced,
		      where $\lpo{nil}$ returns the nil-terminator
		      of an associative symbol and $\lpo{as}$
		      performs type ascription:
		      $$
			      \kw{rule}\ \prn{\lpo{nil}\ \bar s\ \ptrn T}
			      ↪ \prn{\lpo{as}\ \tm{t_\nil}\ \ptrn T}
		      $$

		\item If $δ$ is a macro definition,
		      $δ ⊢ s(\vec ρ) ≔ t$, then
		      $\cmd{δ}$ is
		      $\kw{symbol}\ \bar s\ \ctx{\vec ρ} :
			      \El{\tm{t'}} ≔ \tm{t}$
		      (with the typing omitted if no $t'$
		      is supplied).

		\item If $δ$ is a program declaration with
		      $δ ⊢ s(\vec ρ) : \earr{\plur t n\, t'}$
		      and cases $\ecase{l_1}{r_1}\ldots\ecase{l_m}{r_m}$,
		      then $\cmd{δ}$ declares $\bar s$,
		      where $T ≔ \El{\tm{ \earr{\plur t n\, t'}}}$:
		      %
		      $$\begin{array}[t]{l}
				      \md{sequential}\ \kw{symbol}\ \bar s\
				      \vec ρ^★ : T
				      \mtt{;}
				      \\[1mm]
				      \begin{array}[t]{@{}l@{\ }l@{\ {↪}\ }l}
					      \kw{rule} & \tm{l_1}^{\ptrn{}} & \tm{r_1}^{\ptrn{}}
					      \\ \cdots \\
					      \kw{with} & \tm{l_m}^{\ptrn{}} & \tm{r_m}^{\ptrn{}}\mtt{;}
				      \end{array}
				      % \quad (j = 1 \ldots m)
			      \end{array}$$
		      %
		      where $\vec ρ^★ ≔ \mbf{impl}(\vec ρ, T)$ binds
		      all of the free parameters in $T$ as implicits.

		\item If $δ$ is a rule declaration with
		      $\attr{args}\ \epar{\plur t n}$ of type
		      $\plur τ n$, then
		      $δ ⊢ s(\vec ρ) :
			      \earr{
				      \eoquote{t_1}\,\ldots\,\eoquote{t_n}
				      \ τ^★}$,
		      where $τ^★$ encodes the premises and
		      conclusion (\autoref{sec:eo-sig}).
		      An auxiliary symbol $s^★$ eliminates the
		      $\lpo{quote}$ types by rewriting to
		      $\tm{τ^★}$ on terms matching the argument
		      patterns; $s$ then takes the arguments
		      explicitly and uses $s^★$ for its return type:
		      $$
			      \begin{array}[t]{l}
				      \kw{symbol}\ \bar{s}^★\
				      \vec σ : \tm{τ_1} → \ldots → \tm{τ_n} → \TYPE \mtt{;}
				      \\[1mm]
				      \kw{rule}\
				      \bar s^★\,\tm{t_1}^{\ptrn{}}\,\ldots\,\tm{t_n}^{\ptrn{}}
				      ↪ \tm{τ^★}^{\ptrn{}}\mtt{;}
				      \\[1mm]
				      \md{constant}\ \kw{symbol}\ \bar s\
				      (α_1 : \tm{τ_1}) \ldots (α_n : \tm{τ_n}) :
				      \prn{\app{\app{{\app{\bar s^★}{α_1}}}{\ldots}}{α_n}}
			      \end{array}
		      $$
	\end{enumerate}
\end{definition}

\begin{example}\label{ex:cmd}
	We illustrate $\cmd{δ}$ on declarations from
	\autoref{ex:eo-rules}. The polymorphic equality
	$\ev{=}$ and the associative $\ev{or}$ translate by
	case~1, the latter also producing a nil-terminator
	rewrite rule:
	\begin{lpbox}
		\lstinputlisting[style=lp, aboveskip=0pt, belowskip=0pt]{listings/cmd-constants.lp}%
	\end{lpbox}
	%
	$\ev{modus\_ponens}$ has no $\attr{args}$, so case~1
	yields its derived type directly:
	\begin{lpbox}
		\lstinputlisting[style=lp, aboveskip=0pt, belowskip=0pt]{listings/cmd-modus-ponens.lp}%
	\end{lpbox}
	%
	$\ev{cnf\_implies\_pos}$ has the compound argument
	$\attr{args}\epar{\eapp{\ev{=>}}{\ev{F1}\ \ev{F2}}}$,
	triggering case~4: an auxiliary symbol matches the
	argument pattern and rewrites to $\lpo{Proof}\ φ$
	for $φ ≔ \tm{\eapp{\ev{or}}
			{\eapp{\ev{not}}{\eapp{\ev{=>}}{\ev{F1}\ \ev{F2}}}
				\ \eapp{\ev{not}}{\ev{F1}}\ \ev{F2}}}$.
	\begin{lpbox}
		\lstinputlisting[style=lp, aboveskip=0pt, belowskip=0pt]{listings/cmd-rule-aux.lp}%
	\end{lpbox}
\end{example}

% ----------------------------------------------------------
\subsubsection{Translating Proofs.}
% ----------------------------------------------------------

\begin{definition}\label{def:cmd-prf}
	%
	Let $π$ be a proof command. Then $\cmd{π}$
	is the list of LambdaPi commands such that:
	%
	\begin{enumerate}\itemsep6pt
		\item If $π$ is an assumption,
		      then $π ⊢ s : \eapp{\mtt{Proof}}{φ}$:
		      %
		      $$\md{constant}\ \kw{symbol}\ \bar s :
			      \lpo{Proof}\ \tm{φ}\mtt{;}$$

		\item If $π$ is a step,
		      then $π ⊢ s : \eapp{\mtt{Proof}}{φ}$
		      and $π ⊢ s ≔ \eapp{r}{\plur t m\ \plur p n}$:
		      %
		      \begin{align*}
			       & \kw{symbol}\ \bar s
			      ≔ \bar r\ \tm{t_1}\ \ldots\ \tm{t_m}\
			      \bar p_1\ \ldots\ \bar p_n\mtt{;}
			      \\
			       & \kw{assert}\ ⊢\ \bar s :
			      \lpo{Proof}\ \tm{φ}\mtt{;}
		      \end{align*}
		      %
		      The body applies rule $\bar r$ to its
		      arguments then premises.
		      %
		      The type is checked separately via
		      $\kw{assert}$, which uses LambdaPi's
		      type inference to fill implicit arguments.
	\end{enumerate}
\end{definition}

\begin{example}\label{ex:proof}
	A \texttt{QF\_UF} proof script (top) deriving
	$\false$ from $\eapp{\ev =}{\ev{x0}\ \ev{x0}}$
	via $\ev{eq\_refl}$ and $\ev{eq\_resolve}$, and
	its LambdaPi translation (bottom) generated by
	\autoref{def:cmd} and \autoref{def:cmd-prf}:

	\eofile{listings/proof.eo}%
	\vspace{-\medskipamount}%
	\lpfile{listings/proof.lp}

	The preamble's $\ev{U}$, $\ev{x0}$ translate by
	case~1 and the $\kw{define}$s by case~2 of
	\autoref{def:cmd}; the assumption $\ev{@p1}$
	becomes a $\md{constant}$ symbol of type
	$\lpo{Proof}\ \tm{φ}$ (case~1 of
	\autoref{def:cmd-prf}); each step applies its
	rule to arguments then premises, with $\kw{assert}$
	verifying the conclusion via type inference.
\end{example}

\end{document}
